\chapter{Trabalhos futuros}
%não se esqueça de definir uma label única para utilizar no comando \ref
\label{chap:trabalhos_futuros}

Chegamos à parte onde você admite, com classe e um leve toque de vergonha, que não deu tempo de fazer tudo o que queria... Essa seção é basicamente o “prometo que da próxima vez eu entrego direitinho, fessor”.

Os trabalhos futuros são aquela lista de tarefas que o pesquisador deixa para alguém do futuro (geralmente ele mesmo, quando for fazer o mestrado xD). Aqui é onde você joga no ar ideias mirabolantes, que no final das contas se resumem a um “Seria legal fazer isso, mas não deu tempo.”

Também é pode sonhar alto: “No futuro, pretende-se integrar este sistema com uma IA baseada em aprendizado profundo, IoT e blockchain.” Tudo isso partindo de um projeto que mal roda no seu notebook de 8 anos atrás com 4GB de RAM.

O importante aqui é mostrar que você tem visão, ambição... e que sabe que seu trabalho é só o primeiro passo de uma longa estrada, possivelmente cheia de melhorias sugeridas pela banca.

Enfim, essa seção é onde você tenta plantar a sementinha da dúvida nos avaliadores: “Será que ele não fez porque era difícil ou porque foi estratégico deixar pro futuro?” Spoiler: foi porque era difícil mesmo.
